% !TEX root = ../notes.tex

\documentclass[../notes.tex]{subfiles}

\begin{document}

\section{March 2}
There will be a quiz on Wednesday on some hopefully basic concepts.

\subsection{The Modular Character}
Let's explain how induction works in general.
\begin{definition}[induction]
	Let $H$ be a closed algebraic subgroup of an algebraic group $G$ over a nonarchimedean local field $F$. For a representation $\tau$ of $H$, we define the \textit{induction} $\op{Ind}^G_H\tau$ to be the representation consisting of those functions $f\colon G(F)\to V$ with the following properties.
	\begin{itemize}
		\item Smooth: there is an open compact subgroup $K\subseteq G$ such that $f(gk)=f(g)$ for all $g\in G$ and $k\in K$.
		\item Induction: for all $h\in H(F)$, one has $f(hx)=\delta_{G/H}(h)^{1/2}\tau(h)f(x)$ for all $x\in G$.
	\end{itemize}
	The $G$-action is given by $(gf)(x)\coloneqq f(xg)$.
\end{definition}
Once again, we need to explain the presence of the character $\delta_{G/H}$, which is similar to our discussion from last class.
\begin{remark}
	Top differential forms on the quotient $X\coloneqq H(F)\backslash G(F)$ are given by sections of the line bundle $\omega_X$. The tangent space at the identity is simply $\det(\mf g/\mf h)^*$, so we conclude that $\omega_X=\det(\mf g/\mf h)\times^HG$.
\end{remark}
\begin{notation}
	Given a line bundle $\mc L$ on a $p$-adic manifold $X$ with transition functions $\{g_{ij}\}$, we let $\left|\mc L\right|$ denote the real line bundle, whose transition functions are $\{\left|g_{ij}\right|\}$.
\end{notation}
\begin{definition}[density]
	Fix a $p$-adic manifold $X$. Then a \textit{density} is a section of the real line bundle $\left|\omega_X\right|$.
\end{definition}
\begin{remark}
	For any $p$-adic manifold $X$, one can integrate compactly supported densities, simply by choosing local coordinates and evaluating the integral along some choice of partition of unity. In particular, one can evaluate local integrals against some choice of Haar measure on the base local field $F$.
\end{remark}
We are now ready to define $\delta_{G/H}$.
\begin{definition}
	Let $H$ be a closed algebraic subgroup of an algebraic group $G$ over a local field $F$, and let their Lie algebras be $\mf h$ and $\mf g$, respectively. Then we define $\delta_{G/H}\colon H\to\RR^+$ by
	\[\delta_{G/H}(x)\coloneqq\left|\det\left(\op{Ad}_x;(\mf g/\mf h)^*\right)\right|.\]
\end{definition}
By undoing the dual and then lifting from the quotient, we see that
\[\det\left(\op{Ad}_x;(\mf g/\mf h)^*\right)=\frac{\det(\op{Ad}_x;\mf h)}{\det(\op{Ad}_x;\mf g)}.\]
These individual numerator and denominator characters are interesting in their own right.
\begin{definition}[modular character]
	Fix a locally compact topological group $G$ with left-invariant Haar measure $\mu$. For any $g\in G$, we note that $\op{Ad}_g^*\mu$ continues to be a left-invariant Haar measure, so it is the same as $\mu$ up to a unique scalar. We may thus define the \textit{modular character}
	\[\delta_G(g)\coloneqq\frac{\op{Ad}_g^*\mu}{\mu}.\]
	Note that $\delta_G$ is independent of the choice of $\mu$.
\end{definition}
\begin{remark}
	If $G$ is an affine algebraic group over a local field $F$ with Lie algebra $\mf g$, then one can show that
	\[\delta_G(g)=\left|\det(\op{Ad}_g;\mf g)\right|^{-1}.\]
	Indeed, the left-invariant Haar measure is simply given (algebraically!) by a choice of top form in $\land^n\mf g^*$, where $n$ is the dimension of $G$. Then one sees that the scalar adjusting such a top form (and thus the Haar measure) is given by the adjoint action on $\land^n\mf g^*$, which is what is given above.
\end{remark}
\begin{remark}
	It turns out that $G$ is unimodular if and only if $\delta_G$ is trivial. Indeed, $\delta_G$ is trivial if and only if $\mu$ is conjugation-invariant. But $\mu$ being conjugation-invariant while already being left-invariant is equivalent to being right-invariant and conjugation-invariant. For example, it follows that abelian groups have trivial modular character.
\end{remark}
\begin{example}
	Let $G$ be an algebraic group over a local field $F$. Note that $g\mapsto\det(\op{Ad}_g;\mf g)$ is some character $G\to\mathbb G_m$, so if $G$ admits no such nontrivial homomorphism, then $\delta_G$ is automatically trivial. For example, this holds if $G$ is unipotent (namely, an extension of $\mathbb G_a$s) or semisimple.
\end{example}
\begin{example}
	Because reductive groups are covered by a product of a semisimple group and an abelian torus, it follows that reductive groups are unimodular.
\end{example}
\begin{example}
	Consider the algebraic group $B\subseteq\op{GL}_2$ of upper-triangular matrices over a local field $F$. Then $\begin{bsmallmatrix}
		a & b \\ 0 & d
	\end{bsmallmatrix}$ acts on the Lie algebra $\mf b$ by
	\[\begin{bmatrix}
		x & y \\ 0 & z
	\end{bmatrix}\mapsto\begin{bmatrix}
		x & (a/d)y \\ 0 & z
	\end{bmatrix},\]
	so it follows that $\delta_B\left(\begin{bsmallmatrix}
		a & b \\ 0 & d
	\end{bsmallmatrix}\right)=\left|a/d\right|$.
\end{example}

\subsection{Back to Induction}
We are still owed an explanation why we would want to include this character $\delta_{G/H}^{1/2}$ in our definition of induction. Here is such an application.
\begin{proposition}
	Let $P\subseteq G$ be a parabolic subgroup of a reductive group $G$ over a local field $F$. If $\tau$ and $\tau^\lor$ are dual admissible representations of $P(F)$, then $\op{Ind}_P^G\tau$ and $\op{Ind}_P^G\tau^\lor$ are dual admissible representations.
\end{proposition}
\begin{proof}
	We will merely exhibit the pairing between these representations. Note that $P\backslash G$ is projective, and $P(F)\backslash G(F)$ is compact. The moral is that our integration theory at most allows us to integrate sections of $\left|\omega_{P(F)\backslash G(F)}\right|$. Here, note that $\left|\omega_{P(F)\backslash G(F)}\right|$ is an equivariant real line bundle $\RR\times^PG$, where $P$ acts on $R$ by $\delta_{G/P}$, which in this case agrees with $\delta_P$.

	Now, let $\tau$ and $\tau^\lor$ be dual representations of $P$. Then we may define a pairing
	\[\op{Ind}_P^G\tau\times\op{Ind}_P^G\tau^\lor\to\CC\]
	by
	\[\langle f,f'\rangle\coloneqq\int_{P(F)\backslash G(F)}\langle f(g),f'(g)\rangle\,dg.\]
	Namely, for this pairing to make sense, we need to be able to realize $\varphi\colon g\mapsto\langle f(g),f(g')\rangle$ as a density on $P(F)\backslash G(F)$. Well, for any $p\in P(F)$, we then see that
	\[\langle f(pg),f'(pg)\rangle=\delta_P(p)\langle f(g),f'(g)\rangle,\]
	so $\varphi$ is a section of $\left|\omega_{P(F)\backslash G(F)}\right|$ on $P\backslash G$.

	Of course, it remains to check that we have actually provided a duality. Well, once one shows that $\op{Ind}_P^G\tau$ is actually an admissible representation, one can pass to $K$-finite vectors everywhere (for some open compact subgroup $K\subseteq G(F)$) to check that this is a perfect pairing on finite-dimensional vector spaces.
\end{proof}
\begin{remark}
	We have slid the fact that $\op{Ind}_P^G\tau$ is admissible under the rug because it is on the homework. However, this argument crucially uses that $P$ is parabolic: one does not expect general inductions to send admissible representations to admissible representations!
\end{remark}
Anyway, let's see some examples.
\begin{example}
	Note that $\op{Ind}_P^G\delta_P^{-1/2}$ contains the constant functions. Indeed, by definition, functions in $\op{Ind}_P^G\delta_P^{-1/2}$ are simply smooth functions on $P\backslash G$.
\end{example}
\begin{example}
	Dually, we see that $\op{Ind}_P^G\delta_P^{1/2}$ consists of sections of $\left|\omega_{P\backslash G}\right|$. Thus, there is a quotient map to the trivial representation $\CC$ given by integration against $\int_{P(F)\backslash G(F)}$.
\end{example}
Let's say something about the representation theory of $\op{GL}_2$.
\begin{example}
	Parabolic induction produces many irreducible representations of $\op{GL}_2(F)$. Indeed, let $B\subseteq\op{GL}_2$ be the upper-triangular subgroup, and let $L$ be the diagonal subgroup, which is the Levi quotient. Then $L(F)=F^\times\times F^\times$, so representations of $L$ are given by a pair of characters $\chi_1,\chi_2\colon F^\times\to\CC^\times$. With this description, one can calculate that $\delta_B$ descends to the character $\left(\left|\cdot\right|,\left|\cdot\right|^{-1}\right)$ on $L(F)$. By doing some Mackey theory, it turns out that $\op{Ind}_P^G(\chi_1,\chi_2)$ is irreducible except in the following situations.
	\begin{itemize}
		\item If $\chi_1/\chi_2=\left|\cdot\right|$, then $\op{Ind}_P^G(\chi_1,\chi_2)$ admits a one-dimensional quotient (and has irreducible kernel).
		\item If $\chi_2/\chi_1=\left|\cdot\right|$, then $\op{Ind}_P^G(\chi_1,\chi_2)$ admits a one-dimensional subrepresentation (and has irreducible quotient).
	\end{itemize}
	By finding some intertwining operators, one can show that $\op{Ind}_B^G(\chi_1,\chi_2)\cong\op{Ind}_B^G(\chi_2,\chi_1)$.
\end{example}
\begin{remark}
	The moral is that one roughly has a map from $(F^\times\times F^\times)/S_2$ to irreducible representations of $\op{GL}_2(F)$. It turns out that there are more representations, which are the ``cuspidal'' ones. The other ones are parameterized by characters of $E^\times$ (up to an $S_2$-symmetry), where $E$ is a quadratic extension of $F$. One way to construct these representations is by the Weil representation; another way is to do (compact) induction from general compact open subgroups.
\end{remark}
\begin{remark}[base change]
	Let $\sigma$ be the automorphism of a quadratic extension $E$ of a local field $F$. Given a smooth irreducible representation $\pi$ of $\op{GL}_2(E)$, one may hope that $\pi\cong\pi^\sigma$ would imply that $\pi$ ``comes from'' $\op{GL}_2(F)$, but there is no obvious way to make this association. This is one instance of Langlands's functoriality principle. In the modern era, we are able to do this descent and even give the character of the descended representation, but we have no direct construction.
\end{remark}

\subsection{Spherical Representations}
For most calculations, we will be interested in certain subclasses of representations.
\begin{definition}[spherical]
	Fix an affine algebraic group $G$ over $\OO_F$. Then an irreducible smooth representation of $G(F)$ is \textit{spherical} if and only if $V^{G(\OO_F)}\ne0$.
\end{definition}
\begin{theorem}
	The spherical representations of $\op{GL}_n(F)$ are in bijection with $n$-tuples in $\left(\CC^\times\right)^n$ up to permutation.
\end{theorem}
To begin, we will have to construct some spherical representations. For the rest of this subsection, $G$ is $\op{GL}_n$, and $B\subseteq G$ is the subgroup of upper-triangular matrices. Then $B$ has Levi quotient $L$ given by the diagonal subgroup of $G$. Now, given any $n$ characters $(\chi_1,\ldots,\chi_n)$ of $F^\times$, we can construct the parabolic induction
\[\op{Ind}_B^G(\chi_1,\ldots,\chi_n).\]
Suppose further that each $\chi_i$ is unramified, in the sense that it factors through $\left|\cdot\right|$. In particular, $\chi_i$ is uniquely determined by the value $\chi_i(\varpi)\in\CC^\times$, where $\varpi\in\OO_F$ is any choice of uniformizer. Then our $n$-tuple of characters $(\chi_1,\ldots,\chi_n)$ is basically given by the vector $(\chi_1(\varpi),\ldots,\chi_n(\varpi))\in\left(\CC^\times\right)^n$.

Let's check that the constructed representation in fact admits a vector fixed by $K\coloneqq\op{GL}_n(\OO_F)$. Well, we are on the hunt for functions $f\colon G\to\CC$ such that
\[f(bxk)=\delta_B(b)^{1/2}(\chi_1(b)\cdots\chi_n(b))f(x)\]
for any $b\in B$ and $k\in K$. By doing some explicit calculations, one has the following.
\begin{proposition}[Iwasawa decomposition]
	Continue with $G\coloneqq\op{GL}_n(F)$. Let $U\subseteq G$ be the subgroup of strictly upper-triangular matrices. Then there is a decomposition
	\[\op{GL}_n(F)=\bigsqcup_{\lambda\in\ZZ^n}U_n(F)\op{diag}\left(\varpi^{\lambda_1},\ldots,\varpi^{\lambda_n}\right)\op{GL}_n(\OO_F).\]
\end{proposition}
\begin{proof}[Sketch]
	We would like to show that we can modify any $g\in\op{GL}_n(F)$ by an element of $U(F)$ on the left and by $K$ on the right to make it in the form $\op{diag}\left(\varpi^{\lambda_1},\ldots,\varpi^{\lambda_n}\right)$. This means that we are allowed column operations which are invertible over $\OO$ and row operations which preserve row ordering (without scaling). Then one does some induction with a version of the Gram--Schmidt process.
\end{proof}
\begin{remark}
	Over $\RR$, one has an analogous decomposition
	\[\op{GL}_n(\RR)=U(\RR)\left(\RR^+\right)^n\op O(n).\]
	This is more or less a Gram--Schmidt reduction to turn any basis of $\op{GL}_n(\RR)$ into an orthonormal one.
\end{remark}
The moral is that the function $f$ that we are hunting for is uniquely determined by its value at $1$, so there is at most one such function $f$. After doing a little unwinding, it remains to check that having $bk=1$ is consistent with $f(bk)=f(1)$. But having $b\in B(F)$ and $k\in K$ with $bk=1$, we see that $b,k\in B\cap K$, so $b$ and $k$ are upper-triangular with diagonal entries contained in $\OO^\times$. But $\delta_B$ and the $\chi_i$s all vanish on $\OO^\times$ by construction!
\begin{remark}
	In fact, this argument shows that the only parabolic inductions which are unramified must come from unramified characters.
\end{remark}

\end{document}